\documentclass[aspectratio=169,12pt]{beamer}
%	\usetheme{Antibes}
	\usetheme{Frankfurt}

%\usepackage{amsmath}

\usepackage{graphics}
\graphicspath{{../img}}

%\definecolor{VerdeApostila}{rgb}{0.058823529411764705, 0.6862745098039216, 0.30980392156862746}
%\usecolortheme[named=VerdeApostila]{structure}

\newcommand{\modo}[2]{
%	#2
} % claro ou escuro
\newcommand{\quebra}{
	\break
}

\newcommand{\iniciar}[3]{

	\begin{frame}[plain]
		\maketitle
		\centering
		
		\textbf{Aula #1}: 
		#2 
		 
		#3
	  \end{frame}

%	\section*{Nesta aula}
	\begin{frame}{Nesta aula}
		\begin{multicols}{3}
			\tableofcontents
		\end{multicols}	
	  \end{frame}  

}

\newcommand{\terminar}[1]{

\section{Exercícios #1}
\begin{frame}[allowframebreaks]{Exercícios}	% containsverbatim
	\input{aula#1_exercícios}	
  \end{frame}	
  
\section*{}  

\begin{frame}[plain]
	\maketitle
	
	\centering	
	\Large
	\textbf{Dúvidas?}
		
	\url{akira.demenech@uel.br}
  \end{frame}

\end{document}

}

%\usepackage[brazil]{babel}	
%\usepackage[T1]{fontenc}

\usepackage{xcolor}

\usepackage{amsmath}
\usepackage{mathdots}

\usepackage{multicol}

\usepackage{hyperref}
\newcommand{\link}[2]{\href{#1}{\texttt{#2}}}

\title{Minicurso de Python}
\author{Guilherme Akira Demenech Mori}
\date{}

\institute{
%	\link{https://www.facebook.com/codificajr/}{
%		\includegraphics[width=0.3\textwidth]{\modo{Logo pronta.pdf}{logo clara.pdf}}
%	}
}

\begin{document}
	\modo{}{\color{white}}	

\iniciar{4}{Arquivos}{TXT e CSV}		  



\section{Texto}

\begin{frame}[containsverbatim]{Leitura}
%	conteúdo...

	Ler as linhas de um arquivo:

	\begin{verbatim}
		with open('nome_do_arquivo.txt', 'r') as leitor:			
			for linha in leitor:				
				print(linha)
	\end{verbatim}
	
	\vfill
	
	Caso o arquivo seja codificado em UTF-8:
	
	\begin{verbatim}
		with open('nome_do_arquivo.txt', 'r', encoding='utf-8') as leitor:			
			for linha in leitor:				
				print(linha)
	\end{verbatim}
\end{frame}

\begin{frame}[containsverbatim]{Escrita}
	%	conteúdo...
	
	Escrever em um arquivo:
	
	\begin{verbatim}
		with open('nome_do_arquivo.txt', 'w') as escritor:			
			print(linha, file=escritor)
	\end{verbatim}
	
	\vfill
	
	Para codificar em UTF-8:
	
	\begin{verbatim}
		with open('nome_do_arquivo.txt', 'w', encoding='utf-8') as escritor:									
			print(linha, file=escritor)
	\end{verbatim}
\end{frame}

\section{CSV}

\begin{frame}[containsverbatim]{Leitura}
	%	conteúdo...
	
	\texttt{import csv}
	
	Ler as linhas e colunas de um arquivo CSV:
	
	\begin{verbatim}		
		
		with open('nome_do_arquivo.csv', 'r') as leitor:			
			leitor_csv = csv.reader(leitor, delimiter=';')  
			for linha in leitor_csv:
				for coluna in linha:						
					print(coluna)
	\end{verbatim}
	
	As linhas são retornadas como listas de colunas.
	
	Altere a codificação do leitor do arquivo ou o delimitador do leitor de CSV se for necessário.
	
\end{frame}

\begin{frame}[containsverbatim]{Escrita}
	%	conteúdo...
	
%	\texttt{import csv}
	
	Escrever em um arquivo CSV
	
	\begin{verbatim}				
		with open('nome_do_arquivo.csv', 'w', newline='') as escritor:			
			escritor_csv = csv.writer(escritor, delimiter=';')  
			escritor_csv.writerow(linha)
	\end{verbatim}
	
	As linhas são listas de colunas.
\end{frame}



\terminar{4}
